\section{Conclusion}
This paper adapts the pretrained visual-language models in natural domain to medical AD, with cross-domain generalizability among various modalities and anatomical regions. 
The adaptation involves not only from natural domain to medical domain, but also from high-level semantics to pixel-level segmentation. 
To achieve such goals, a collaborative multi-level feature adaptation method is introduced, where each adaptation is guided by the corresponding visual-language alignments, facilitating segmenting anomalies of diverse forms from medical images.
Coupled with a comparison-based AD strategy, the method enables flexible adaptation for datasets with substantial modality and distribution differences. 
The proposed method outperforms existing methods on zero-/few-shot AC and AS tasks, indicating promising research avenues for future exploration.

\noindent\textbf{Acknowledgment.} This work is supported by the National Key R\&D Program of China (No. 2022ZD0160703), STCSM (No. 22511106101, No. 18DZ2270700, No. 21DZ1100100), 111 plan (No. BP0719010), State Key Laboratory of UHD Video and Audio Production and Presentation,
and the Ministry of Education, Singapore, under its Academic Research Fund Tier 2 (Award Number: MOE-T2EP20122-0006).